%% !!!!POZOR!!!!
% je zapnuté oneside dokument v sablona-prace.tex
% a electronic v nastaveni.tex 
% před odevzdáním bude nejspíš nutno změnit
% a hlavně před případným tiskem

\chapter{Teoretická část studentské práce}

\section{Metody detekce kolizí pohybu}

\subsection{Vzdálenost koncových bodů}

\subsection{Vzdálenost koncových bodů s vektorem pohybu}

\subsection{Rozdělení prostoru na kvadráty}

\newpage

\section{Simulované prostředí}
    Důležitou částí semestrální práce je vytvoření simulovaného
    prostředí pracoviště robotů.

    \subsection{Programy pro vytvoření simulace}
        Jsou tři hlavní programy v kterých 
        lze simulace vytvořit. Program Scene Viewer, industrialPhysics
        a Unity 3D.

        \subsubsection{Scene Viewer}
            Program Scene Viewer je vytvořený přímo od firmy B\&R, tedy 
            patří do jednoho kompletního řešení podporovaného přímo
            vývojáři a nejlépe zdokumentovaným propojením s virtuálním
            i reálným PLC od této firmy.
    
            Největší výhodou je existence vytvořených modelů robotů,
            takže není nutné zabývat se vytvářením vlastního 3D modelu. 
            Simulátor komunikuje s PLC pomocí OPC UA, nebo alternativně
            pomocí PVI.

            Scene Viewer lze dostat zdarma na stránkách B\&R. Při prvním
            spuštění je třeba vytvořit certifikát, který se uloží
            na počítači.

    \subsubsection{industrialPhysics}

    \subsubsection{Unity 3D}